\subsection{\texorpdfstring{同态 \(\operatorname{Tor}^{\mathbf{B}}(\mathbf{P},\mathbf{N})\otimes_{\mathbf{A}}\mathbf{Q}\to \operatorname{Tor}^{\mathbf{B}}(\mathbf{P},\mathbf{N}\otimes_{\mathbf{A}}\mathbf{Q})\) 与 \(\operatorname{Ext}_{\mathbf{B}}(\mathbf{M},\mathbf{N})\otimes_{\mathbf{A}}\mathbf{Q}\rightarrow \operatorname{Ext}_{\mathbf{B}}(\mathbf{M},\mathbf{N}\otimes_{\mathbf{A}}\mathbf{Q})\)}{张量积下的 Tor 与 Ext 同态}}

设 B 是环，N 是 (B, A)-双模，M 是左 B-模，P 是右 B-模，Q 是左 A-模。

根据 X，第 62 页，我们有同态

\[
\gamma_ {1}: \mathrm{H} (\mathrm{L} (\mathrm{P}) \otimes_ {\mathrm{B}} \mathrm{N}) \otimes_ {\mathrm{A}} \mathrm{Q} \rightarrow \mathrm{H} (\mathrm{L} (\mathrm{P}) \otimes_ {\mathrm{B}} \mathrm{N} \otimes_ {\mathrm{A}} \mathrm{Q});
\]

此外（X，第 69 页，命题 5），我们有同构

\[
\begin{array}{c} \psi_ {1} (N): \operatorname{Tor} ^ {B} (P, N) \otimes_ {A} Q \qquad H (L (P) \otimes_ {B} N) \otimes_ {A} Q, \\ \psi_ {1} (N \otimes_ {A} Q): \operatorname{Tor} ^ {B} (P, N \otimes_ {A} Q) \to H (L (P) \otimes_ {B} N \otimes_ {A} Q). \end{array}
\]

典范的 0 次分次同态

\[
\operatorname{Tor} ^ {\mathrm{B}} (\mathrm{P}, \mathrm{N}) \otimes_ {\mathrm{A}} \mathrm{Q} \rightarrow \operatorname{Tor} ^ {\mathrm{B}} (\mathrm{P}, \mathrm{N} \otimes_ {\mathrm{A}} \mathrm{Q})\tag{9}
\]

定义为复合 \(\psi_{\mathrm{P}}(\mathbf{N}\otimes_{\mathbf{A}}\mathbf{Q})^{-1}\circ \gamma_{1}\circ (\psi_{\mathrm{P}}(\mathbf{N})\otimes 1_{\mathbf{Q}})\) .

类似地，由复形的典范态射

\[
\alpha : \operatorname{Homgr} _ {B} (L (M), N) \otimes_ {A} Q \rightarrow \operatorname{Homgr} _ {B} (L (M), N \otimes_ {A} Q)
\]

由此推出同态 H(α)；我们有典范同态（X，第 62 页）

\[
\gamma_ {2}: \mathrm{H} (\operatorname{Homgr} _ {\mathrm{B}} (L (M), N)) \otimes_ {A} Q \rightarrow \mathrm{H} (\operatorname{Homgr} _ {\mathrm{B}} (L (M), N) \otimes_ {A} Q),
\]

和同构（X，第 88 页，命题 5）

\[
\begin{array}{c} \varphi_ {\mathrm{M}} (\mathrm{N}): \mathrm{H} \big (\operatorname{Homgr} _ {\mathrm{B}} (\mathrm{L} (\mathrm{M}), \mathrm{N}) \big) \otimes_ {\mathrm{A}} \mathrm{Q} \to \operatorname{Ext} _ {\mathrm{B}} (\mathrm{M}, \mathrm{N}) \otimes_ {\mathrm{A}} \mathrm{Q}, \\ \varphi_ {\mathrm{M}} (\mathrm{N} \otimes_ {\mathrm{A}} \mathrm{Q}): \mathrm{H} \big (\operatorname{Homgr} _ {\mathrm{B}} (\mathrm{L} (\mathrm{M}), \mathrm{N} \otimes_ {\mathrm{A}} \mathrm{Q}) \big) \to \operatorname{Ext} _ {\mathrm{B}} (\mathrm{M}, \mathrm{N} \otimes_ {\mathrm{A}} \mathrm{Q}). \end{array}
\]

典范的 0 次分次同态

\[
\operatorname{Ext} _ {\mathbf {B}} (\mathbf {M}, \mathbf {N}) \otimes_ {\mathbf {A}} \mathbf {Q} \rightarrow \operatorname{Ext} _ {\mathbf {B}} (\mathbf {M}, \mathbf {N} \otimes_ {\mathbf {A}} \mathbf {Q})\tag{10}
\]

定义为复合

\[
\varphi_ {\mathrm{M}} (\mathrm{N} \otimes_ {\mathrm{A}} \mathrm{Q}) \circ \mathrm{H} (\alpha) \circ \gamma_ {2} \circ \left(\varphi_ {\mathrm{M}} (\mathrm{N}) \otimes 1 _ {\mathrm{Q}}\right) ^ {- 1}.
\]

命题 7. --- a) 若 A-模 Q 是平坦的，则同态 (9) 是双射的。b) 若 A-模 Q 是有限生成投射的，则同态 (10) 是双射的。

\begin{enumerate}
\def\labelenumi{\alph{enumi})}
\setcounter{enumi}{2}
\item
  若 A-模 Q 是平坦的，环 B 是诺特的，且 B-模 M 是有限生成的，则同态 (10) 是双射的。
\item
  若 Q 是平坦的， \(\gamma_{1}\) 是双射的（X，第 66 页，推论 2）。
\item
  若 Q 是有限生成投射的，则它是平坦的，因此 \(\gamma_{2}\) 是双射的，而且 \(\alpha\) 是双射的（II，第 75 页，命题 2，a)）。
\item
  在 c) 的假设下， \(\gamma_{2}\) 是双射的，因为 Q 是平坦的。此外（X，第 53 页，命题 6），存在 M 的分解 L，使得每个 \(\mathbf{L}_{n}\) 都是有限生成自由的；设 \(u: \mathbf{L}(\mathbf{M}) \to \mathbf{L}\) 是同伦等价（X，第 49 页，命题 3 的推论）；在交换图中，
\end{enumerate}

\[
\begin{array}{c c c} \operatorname{Homgr} _ {\mathrm{B}} (\mathrm{L} (\mathrm{M}), \mathrm{N}) \otimes_ {\mathrm{A}} \mathrm{Q} & \xrightarrow {\alpha} & \operatorname{Homgr} _ {\mathrm{B}} (\mathrm{L} (\mathrm{M}), \mathrm{N} \otimes_ {\mathrm{A}} \mathrm{Q}) \\ \uparrow & & \uparrow \\ \operatorname{Homgr} _ {\mathrm{B}} (\mathrm{L}, \mathrm{N}) \otimes_ {\mathrm{A}} \mathrm{Q} & \xrightarrow {\bar {\alpha}} & \operatorname{Homgr} _ {\mathrm{B}} (\mathrm{L}, \mathrm{N} \otimes_ {\mathrm{A}} \mathrm{Q}). \end{array}
\]

由 \(u\) 推出的竖直箭头是同伦等价（X，第 64 页，命题 3 和第 83 页，命题 3），且 \(\overline{\alpha}\) 是双射的（II，第 75 页，命题 2 (ii)）；因此 H(α) 是双射的，且同态 (10) 是双射的。

